\chapter{Why This Text Exists}
\label{ch:overview}

The original manuscript established a useful thesis: people often fail to
initiate beneficial action not because they lack values, but because they are
embedded in a current behavioural frame whose reward structure, cognitive
momentum, and environment all resist transition. That thesis is still the
center of the project. What changes in this draft is the teaching ambition.

This version is written as a textbook prototype rather than a short conceptual
essay. A textbook needs more than correct intuition. It needs ordered
definitions, stable notation, chapter progression, worked cases, and an
exercise layer that lets the reader test whether the framework actually
explains behaviour rather than merely sounding insightful.

\section{Learning Objectives}

After reading this draft, a reader should be able to:
\begin{enumerate}[nosep]
  \item describe a behavioural state as a structured object rather than a mood;
  \item distinguish reward, switching cost, and environmental reinforcement;
  \item identify decision windows at which free will has unusually high
        leverage;
  \item explain why planning changes the geometry of later action rather than
        merely adding motivation;
  \item relate behavioural transitions to empirical work on conscious-state
        stability and state change;
  \item explain why frameworks can emerge from repeated data, coupling, and
        interference rather than appearing as static labels;
  \item identify which advanced dynamical-systems tools clarify the framework
        and which belong only in appendices or frontier discussion;
  \item classify real-life failures using a reusable case vocabulary;
  \item design interventions that lower activation barriers before failure
        occurs.
\end{enumerate}

\section{Book Architecture}

The manuscript is organized around nine teaching jobs:
\begin{enumerate}[nosep]
  \item Chapters~\ref{ch:formal-language} and \ref{ch:dynamics} define the
        language of the model.
  \item Chapter~\ref{ch:planning} explains why planning is structurally
        necessary.
  \item Chapter~\ref{ch:consciousness-biology} introduces the empirical
        consciousness layer and ties state transitions to biological evidence.
  \item Chapter~\ref{ch:emergence} explains how durable frameworks arise from
        repeated interaction, not from a single isolated choice.
  \item Chapter~\ref{ch:advanced-dynamics} upgrades the mathematical language
        from a lightweight barrier model to a richer state-space view.
  \item Chapter~\ref{ch:cases} makes the model concrete across multiple
        domains.
  \item Chapter~\ref{ch:interventions} turns analysis into design patterns.
  \item Chapter~\ref{ch:exercises} tests whether the reader can reuse the
        framework.
  \item The appendices preserve notation, templates, and clearly isolated
        frontier hypotheses.
\end{enumerate}

\section{Scope and Limits}

This is still a semi-formal project. The purpose is to create a strong bridge
between mathematical style and behavioural usefulness. Therefore:
\begin{itemize}[nosep]
  \item the notation is disciplined but intentionally lightweight;
  \item the propositions are explanatory tools, not claims of finished
        empirical law;
  \item empirical neuroscience, theoretical systems language, and speculative
        consciousness-physics proposals are kept in separate evidence lanes;
  \item the case library is designed to improve transfer, not to replace
        psychological literature;
  \item the intervention patterns are production templates for later expansion.
\end{itemize}

\section{Evidence Policy}

This research-backed version uses four lanes:
\begin{enumerate}[nosep]
  \item \textbf{Strong empirical core:} material suitable for the main
        consciousness and state-transition argument.
  \item \textbf{Mechanism-building theory:} emergence and dynamics sources that
        deepen explanation without pretending to be direct measurement.
  \item \textbf{Speculative frontier:} high-risk proposals that remain outside
        the empirical core.
  \item \textbf{Recommended reading:} classic textbooks and major reviews that
        help the reader go deeper chapter by chapter.
\end{enumerate}

\begin{readingbox}
\textbf{Foundation}
\begin{itemize}[nosep]
  \item Stanislas Dehaene, \emph{Consciousness and the Brain}.
  \item Christof Koch, \emph{Consciousness}.
\end{itemize}
\textbf{Modern Research}
\begin{itemize}[nosep]
  \item Use the graded literature map for exact paper intake before future
        chapter expansion.
\end{itemize}
\textbf{Frontier}
\begin{itemize}[nosep]
  \item Frontier readings should be consumed only after the empirical core is
        understood.
\end{itemize}
\end{readingbox}

\begin{summarybox}
This chapter changes the business value of the manuscript. The project is no
longer just a single argument. It is now a reusable teaching system with a
defined learning path, explicit scope, and a structure that can absorb future
proofs, cases, and exercises without collapsing into notes.
\end{summarybox}
